\scanpage{204}
\mspage{60}
Je voudrait généraliser \deleted{\unsure{les}} la notion de flèches $W_0$-biasphériques au cas d'un dérivateur $\mathbb{D}$ quelconque. \deleted{\unsure{Pour}}

\uline{Proposition} 4.4. Soit $\mathbb{D}$ un dérivateur\inserted{sur $\mathrm{Diag}$}, et $u\colon I'\mathbin{\unsure{-|\to}}I$ \unsure{d'une|une} flèche dans $\mathrm{Diag}$. \deleted{Considérons les conditions :} \inserted{Supposons que $v_*$ et $v_!$ existent, et se calculent de façon standard, pour $\forall$ flèche $v$ de $\mathrm{Diag}$. (En fait, il suffit que ce soit le cas pour $u$ et ses localisés $u_{/i}\colon I'_{/i}\mathbin{-}I_{/i}$, ou $u$ et ses colocalisés $\,{}_{i\backslash}u\colon{}_{i\backslash}I'\mathbin{-}{}_{i\backslash}I$.) Considérons les conditions :}

a) $u^*\colon\mathbb{D}(I)\mathbin{-}\mathbb{D}(I')$ pl. fid.

$a')$ $\forall\ F$ dans $\mathbb{D}(I)$, la flèche

\[
 u_!u^*F\longrightarrow F
\]

est un iso. \deleted{(bien supposé que $u_!$ existe)}

$a'')$ $\forall\ F$ dans $\mathbb{D}(I)$, la flèche

\[
 u_*u^{\unsure{u|*}}F\longleftarrow F
\]

est un iso. \deleted{(bien supposé que $u_*$ existe)}

b) $\forall\ i$ \unsure{dans|de} $I$, la flèche $\,{}_{i\backslash}I'\mathbin{\unsure{-|\to}}{}_{i\backslash}I$ est $\mathbb{D}$-asphérique.

$b')$ $\forall\ i$ \unsure{de|dans} $I$, la flèche $I'_{/i}\mathbin{-}I_{/i}$ est $\mathbb{D}$-coasphérique.

$b'')$ $\forall\ \alpha\colon i_0\mathbin{-}\unsure{i_1'|i_1}$ de $I$, la cat $\,{}_{i_0\backslash}\underset{\alpha}{I'}_{/i_1}$ est $\mathbb{D}$-asphérique.

On a alors les implications

\[
 (b\Longleftrightarrow b'\Longleftrightarrow b'')\Longrightarrow
 (a\Longleftrightarrow a'\Longleftrightarrow a'')
\]

\deleted{(où $a',a''$ sont mis sous réserve d'existence de $u_!,u_*$ respectivement), et où la flèche $\Rightarrow$ est valable sous réserve que $u_!$ ou $u_*$ existe. Sous la même réserve,} On a $a\Rightarrow b$ (i.e. toutes les conditions \deleted{\unsure{écrites|ci-dessus}} envisagées sont équivalentes)
% archive sheet number: 204
